\let\oldprintchaptertitle=\printchaptertitle
\renewcommand{\printchaptertitle}[1]{%
	\vspace*{-75pt}
	\oldprintchaptertitle{#1}
}%
\myChapterStar{Abstract}{}{section}
\let\printchaptertitle=\oldprintchaptertitle

Monitoring groundwater quality in the context of per- and polyfluoroalkyl substances (PFAS) contamination raises a major operational challenge. Since the U.S. EPA adopted the first mandatory maximum contaminant levels (MCLs) in 2024 (down to 4~ng/L for PFOA and PFOS under the NPDWR regulation), water system managers must prioritize sampling campaigns without being able to submit every well to costly chromatographic analysis. This need defines the \emph{strict predictive mode}: estimating the risk of regulatory exceedance from geospatial contextual variables alone, with no previously measured PFAS concentration as input. This work addresses two barriers that the existing literature does not tackle simultaneously. The first is a software barrier: geospatial enrichment of raw observations is typically handled by monolithic, non-reusable scripts. We propose \texttt{geoenrich}, a generic engine organized into three layers (an agnostic \emph{contract} layer, a parameterizable \emph{operators} layer, and a \emph{configuration} layer), whose correctness was validated through a parity test against the original pipeline: 77~identical columns, 17~improved, and zero regressions out of 94~columns produced. The second barrier is scientific and stems from the spatial autocorrelation of PFAS contamination, which makes a random train-test split overly optimistic. To address this, all location variables were excluded from model inputs, and five models were trained under an identical protocol on 46,338~California groundwater samples: two tabular baselines (random forest and XGBoost), a heterogeneous graph transformer (HGT) encoding the relational structure between wells and contamination sources, and two hybrid models (Embedding Fusion and Stacking). On the test set, random forest and stacking both reach F1~0.925 (AUC 0.979 and 0.977); XGBoost achieves F1~0.923 (AUC~0.977); the HGT lags behind (F1~0.825, AUC~0.906). Shapley value analysis confirms that hydrogeological context variables and proximity to contamination sources carry most of the predictive signal, validating the removal of coordinates and showing that careful geospatial enrichment alone is sufficient to achieve operational discrimination.

\vspace{1cm}
\noindent\textbf{Keywords:} PFAS, groundwater, geospatial enrichment, machine learning, binary classification, predictive mode, heterogeneous graph transformer, Shapley values, interpretability, California.

\myCleanStarChapterEnd
